% Options for packages loaded elsewhere
\PassOptionsToPackage{unicode}{hyperref}
\PassOptionsToPackage{hyphens}{url}
\documentclass[
]{article}
\usepackage{xcolor}
\usepackage[margin=1in]{geometry}
\usepackage{amsmath,amssymb}
\setcounter{secnumdepth}{-\maxdimen} % remove section numbering
\usepackage{iftex}
\ifPDFTeX
  \usepackage[T1]{fontenc}
  \usepackage[utf8]{inputenc}
  \usepackage{textcomp} % provide euro and other symbols
\else % if luatex or xetex
  \usepackage{unicode-math} % this also loads fontspec
  \defaultfontfeatures{Scale=MatchLowercase}
  \defaultfontfeatures[\rmfamily]{Ligatures=TeX,Scale=1}
\fi
\usepackage{lmodern}
\ifPDFTeX\else
  % xetex/luatex font selection
\fi
% Use upquote if available, for straight quotes in verbatim environments
\IfFileExists{upquote.sty}{\usepackage{upquote}}{}
\IfFileExists{microtype.sty}{% use microtype if available
  \usepackage[]{microtype}
  \UseMicrotypeSet[protrusion]{basicmath} % disable protrusion for tt fonts
}{}
\makeatletter
\@ifundefined{KOMAClassName}{% if non-KOMA class
  \IfFileExists{parskip.sty}{%
    \usepackage{parskip}
  }{% else
    \setlength{\parindent}{0pt}
    \setlength{\parskip}{6pt plus 2pt minus 1pt}}
}{% if KOMA class
  \KOMAoptions{parskip=half}}
\makeatother
\usepackage{longtable,booktabs,array}
\newcounter{none} % for unnumbered tables
\usepackage{calc} % for calculating minipage widths
% Correct order of tables after \paragraph or \subparagraph
\usepackage{etoolbox}
\makeatletter
\patchcmd\longtable{\par}{\if@noskipsec\mbox{}\fi\par}{}{}
\makeatother
% Allow footnotes in longtable head/foot
\IfFileExists{footnotehyper.sty}{\usepackage{footnotehyper}}{\usepackage{footnote}}
\makesavenoteenv{longtable}
\setlength{\emergencystretch}{3em} % prevent overfull lines
\providecommand{\tightlist}{%
  \setlength{\itemsep}{0pt}\setlength{\parskip}{0pt}}
\usepackage{bookmark}
\IfFileExists{xurl.sty}{\usepackage{xurl}}{} % add URL line breaks if available
\urlstyle{same}
\hypersetup{
  hidelinks,
  pdfcreator={LaTeX via pandoc}}

\author{}
\date{}

\begin{document}

\section{Quark and Lepton Mixing Matrices from Lattice Geometry with
Zero Free
Parameters}\label{quark-and-lepton-mixing-matrices-from-lattice-geometry-with-zero-free-parameters}

\textbf{Jonathan D. Wollenberg}

March 5, 2026

\begin{center}\rule{0.5\linewidth}{0.5pt}\end{center}

\subsection{Abstract}\label{abstract}

We present closed-form expressions for all four CKM parameters and all
four PMNS parameters that reproduce the observed quark and lepton mixing
matrices with zero free parameters. The CKM matrix is constructed in the
standard PDG parametrization with angles theta\_12 =
arcsin(sqrt(m\_d/m\_s + m\_u/m\_c)), theta\_23 =
arcsin(sqrt(m\_u/m\_c)), theta\_13 = arcsin(sqrt(m\_u/m\_t)), and CP
phase delta = arccos(5/12), yielding all nine matrix elements within 1.4
sigma of PDG 2024 values (mean error 0.64\%). The PMNS matrix is
obtained by applying a single geometric rotation R(theta, axis) to the
tribimaximal mixing matrix, with theta = arcsin((m\_e/m\_mu)\^{}(1/3))
and a rotation axis determined by lepton-to-proton mass ratios, yielding
all three angles within 1 sigma of NuFIT 6.0 (normal ordering). The
quark sector uses mass ratios raised to the 1/2 power (surface
geometry), while the lepton sector uses the 1/3 power (bulk geometry).
Both CP phases derive from the tetrahedral dihedral angle arccos(+/-1/3)
in d = 3 dimensions, with the CKM phase receiving a boundary correction
to arccos(5/12). The formulas require only measured fermion masses as
input and contain no fitted or adjustable parameters. The geometric
origin of these relations, within the framework of standing-wave modes
on a three-dimensional lattice, is briefly discussed.

\begin{center}\rule{0.5\linewidth}{0.5pt}\end{center}

\subsection{1. Introduction}\label{introduction}

The Standard Model of particle physics contains 19 free parameters, of
which 10 describe the quark and lepton mixing matrices: four CKM
parameters (three angles and one CP phase) and four PMNS parameters
(three angles and one CP phase), plus two neutrino mass-squared
differences. These parameters are determined experimentally but have no
theoretical explanation within the Standard Model itself.

Several empirical relations between mixing angles and fermion mass
ratios have been noted in the literature. The Gatto-Sartori-Tonin
relation V\_us \textasciitilde{} sqrt(m\_d/m\_s) {[}1{]} and its
extension to the full Cabibbo angle via quadrature V\_us =
sqrt(m\_d/m\_s + m\_u/m\_c) {[}2{]} are well-known. The tribimaximal
(TBM) ansatz for neutrino mixing {[}3{]} successfully predicted the
approximate structure of the PMNS matrix before theta\_13 was measured.
However, these relations have generally been treated as separate
observations without a unifying geometric framework, and most require at
least one fitted parameter.

In this paper, we present a complete set of formulas for all eight
mixing parameters (four CKM, four PMNS) that:

\begin{enumerate}
\def\labelenumi{\arabic{enumi}.}
\tightlist
\item
  Contain zero free parameters --- all inputs are measured fermion
  masses and the integer d = 3.
\item
  Reproduce all observed mixing angles within experimental uncertainty.
\item
  Share a common geometric structure: mass ratios raised to a power
  determined by the dimensionality of the coupling boundary.
\end{enumerate}

The key structural observation is that quark mixing angles involve mass
ratios to the 1/2 power, while lepton mixing angles involve mass ratios
to the 1/3 power. We interpret this as the difference between coupling
at a (d-1) = 2 dimensional surface (quarks, confined inside the proton)
and coupling through d = 3 dimensional bulk (leptons, which are free
particles). Both CP-violating phases derive from the dihedral angle of a
regular tetrahedron in three dimensions: arccos(1/3) for quarks (with a
boundary correction) and arccos(-1/3) for leptons.

\begin{center}\rule{0.5\linewidth}{0.5pt}\end{center}

\subsection{2. Framework}\label{framework}

\subsubsection{2.1 Lattice foundation}\label{lattice-foundation}

The formulas presented here arise from a framework called Geometric Wave
Theory (GWT), in which fundamental particles are standing-wave modes
(breathers) of a sine-Gordon field on a d-dimensional cubic lattice with
Planck spacing. The Lagrangian density is:

\begin{verbatim}
L = sum_{<i,j>} [ (1/2)(phi_i - phi_j)^2 + (1/pi^2)(1 - cos(pi * phi_i)) ]
\end{verbatim}

where phi\_i is the displacement field at lattice site i and the sum
runs over nearest neighbors. This lattice has a single structural input:
the spatial dimension d = 3. From its breather spectrum and tunneling
amplitudes, all fermion masses are determined by:

\begin{verbatim}
m(n, p) = (16/pi^2) * sin(n * gamma) * exp(-16p/pi^2) * m_Planck
\end{verbatim}

where gamma = pi/(16\emph{pi - 2) is the breather coupling, n is the
harmonic mode number (an integer from 1 to N = d } 2\^{}d = 24), and p
is the tunneling depth. Both n and p are integers determined by lattice
symmetry.

The nine charged fermion masses predicted by this formula, with specific
(n, p) assignments derived from the lattice harmonic structure, agree
with observation to a mean accuracy of 1.5\% {[}4{]}. These masses serve
as inputs to the mixing angle formulas below. For transparency, we
present results using both GWT-predicted masses and PDG experimental
masses, so that the mixing formulas can be evaluated independently of
the mass predictions.

\subsubsection{2.2 Surface vs.~bulk
geometry}\label{surface-vs.-bulk-geometry}

A central distinction in this framework is between particles that are
confined inside a topological defect (the proton, modeled as a kink on
the lattice) and particles that propagate freely on the full lattice.

\begin{itemize}
\item
  \textbf{Quarks} are confined within the proton. Their weak-interaction
  coupling occurs at the proton's (d-1) = 2 dimensional boundary.
  Overlap integrals between standing-wave modes on a 2D surface yield
  amplitudes proportional to (m\_i/m\_j)\^{}(1/(d-1)) =
  (m\_i/m\_j)\^{}(1/2).
\item
  \textbf{Leptons} are free particles extending through the full d = 3
  dimensional lattice. Their mixing amplitudes scale as
  (m\_i/m\_j)\^{}(1/d) = (m\_i/m\_j)\^{}(1/3).
\end{itemize}

This distinction --- square root for quarks, cube root for leptons ---
is the single structural difference between the CKM and PMNS
derivations.

\subsubsection{2.3 Tetrahedral geometry and CP
phases}\label{tetrahedral-geometry-and-cp-phases}

The regular tetrahedron is the unique symmetric solid in d = 3
dimensions with d + 1 = 4 vertices. Its dihedral angle is arccos(1/d) =
arccos(1/3) = 70.53 degrees. In the lattice framework, the four vertices
correspond to the four possible orientations of a kink defect in the
d-cube, and CP violation arises from the chirality of wave propagation
around these vertices.

The two CP phases are:

\begin{itemize}
\tightlist
\item
  \textbf{PMNS}: delta\_PMNS = arccos(-1/d) = arccos(-1/3) = 109.47
  degrees (bulk, negative handedness)
\item
  \textbf{CKM}: delta\_CKM = arccos((d+2)/(d(d+1))) = arccos(5/12) =
  65.38 degrees (surface-corrected, positive handedness)
\end{itemize}

The CKM correction shifts cos(delta) from 1/3 to 5/12 = 1/3 + 1/12. The
additional 1/(d(d-1)\^{}2) = 1/12 term arises because quark mixing
occurs at the (d-1)-dimensional boundary of the proton, adding a surface
contribution to the bulk dihedral angle.

\begin{center}\rule{0.5\linewidth}{0.5pt}\end{center}

\subsection{3. CKM Matrix}\label{ckm-matrix}

\subsubsection{3.1 Construction}\label{construction}

The CKM matrix is constructed in the standard PDG parametrization:

\begin{verbatim}
V_CKM = R_23(theta_23) * R_13(theta_13, delta) * R_12(theta_12)
\end{verbatim}

with four geometric parameters:

{\def\LTcaptype{none} % do not increment counter
\begin{longtable}[]{@{}lll@{}}
\toprule\noalign{}
Parameter & Formula & Value \\
\midrule\noalign{}
\endhead
\bottomrule\noalign{}
\endlastfoot
theta\_12 & arcsin(sqrt(m\_d/m\_s + m\_u/m\_c)) & 12.957 degrees \\
theta\_23 & arcsin(sqrt(m\_u/m\_c)) & 2.392 degrees \\
theta\_13 & arcsin(sqrt(m\_u/m\_t)) & 0.203 degrees \\
delta & arccos(5/12) & 65.376 degrees \\
\end{longtable}
}

The theta\_12 formula is the well-known quadrature relation: the
down-type and up-type sectors contribute perpendicular rotations in
flavor space, so their squares add: sin\^{}2(theta\_12) = m\_d/m\_s +
m\_u/m\_c.~The theta\_23 formula identifies V\_cb with the ``up-type
Cabibbo angle'' sqrt(m\_u/m\_c). The theta\_13 formula gives V\_ub as
the direct 1-3 surface overlap sqrt(m\_u/m\_t).

\subsubsection{3.2 Results}\label{results}

Using GWT quark masses (m\_u = 2.214, m\_d = 4.783, m\_s = 98.56, m\_c =
1271, m\_b = 4312, m\_t = 176547 MeV):

{\def\LTcaptype{none} % do not increment counter
\begin{longtable}[]{@{}lllll@{}}
\toprule\noalign{}
Element & Predicted & PDG 2024 & Uncertainty & Pull (sigma) \\
\midrule\noalign{}
\endhead
\bottomrule\noalign{}
\endlastfoot
V\_ud & 0.97453 & 0.97435 & 0.00016 & +1.1 \\
V\_us & 0.22422 & 0.22500 & 0.00067 & -1.2 \\
V\_ub & 0.003541 & 0.00369 & 0.00011 & -1.4 \\
V\_cd & 0.22408 & 0.22486 & 0.00067 & -1.2 \\
V\_cs & 0.97368 & 0.97349 & 0.00016 & +1.2 \\
V\_cb & 0.04173 & 0.04182 & 0.00085 & -0.1 \\
V\_td & 0.00852 & 0.00857 & 0.00020 & -0.3 \\
V\_ts & 0.04101 & 0.04110 & 0.00085 & -0.1 \\
V\_tb & 0.99912 & 0.99912 & 0.00004 & +0.1 \\
\end{longtable}
}

Mean \textbar error\textbar: 0.64\%. Maximum pull: 1.4 sigma (V\_ub).
Unitarity is exact by construction.

The Jarlskog invariant, which measures the overall strength of CP
violation:

\begin{verbatim}
J = Im(V_ud * V_cs * V_us* * V_cd*) = 2.93 x 10^-5
\end{verbatim}

Observed: J = (3.08 +/- 0.13) x 10\^{}-5. Deviation: -4.8\%.

\subsubsection{3.3 Comparison with previous
formulas}\label{comparison-with-previous-formulas}

The formula V\_cb = sqrt(m\_u/m\_c) replaces the previous Wolfenstein
parametrization V\_cb = A * lambda\^{}2 with A = sqrt(2/3). Both give
sub-percent accuracy, but the new formula is simpler (one mass ratio
vs.~a composite expression) and eliminates the need for the Wolfenstein
amplitude A as a separate geometric input. The V\_us quadrature formula
and V\_ub = sqrt(m\_u/m\_t) are unchanged from earlier work {[}2, 5{]}.

\begin{center}\rule{0.5\linewidth}{0.5pt}\end{center}

\subsection{4. PMNS Matrix}\label{pmns-matrix}

\subsubsection{4.1 Construction}\label{construction-1}

The PMNS matrix is obtained by applying a single rotation to the
tribimaximal (TBM) base:

\begin{verbatim}
U_PMNS = R(theta_corr, n_hat) * U_TBM
\end{verbatim}

where U\_TBM is the Harrison-Perkins-Scott tribimaximal matrix {[}3{]}:

\begin{verbatim}
U_TBM = | sqrt(2/3)   1/sqrt(3)   0        |
        | -1/sqrt(6)  1/sqrt(3)   1/sqrt(2) |
        | 1/sqrt(6)   -1/sqrt(3)  1/sqrt(2) |
\end{verbatim}

The correction parameters are:

\begin{verbatim}
theta_corr = arcsin((m_e / m_mu)^(1/3)) = 9.74 degrees

n_hat = (-1, sqrt(3), -(m_tau/m_p)^(1/3)) / |(-1, sqrt(3), -(m_tau/m_p)^(1/3))|
\end{verbatim}

The correction angle uses the cube root (1/d = 1/3, bulk geometry) of
the electron-to-muon mass ratio. The rotation axis has components (-1,
sqrt(3)) in the TBM degenerate subspace, with a third component
proportional to the tau-to-proton mass ratio raised to the 1/3 power.
This wrapping factor (m\_tau/m\_p)\^{}(1/3) accounts for the tau
lepton's standing-wave extent relative to the proton's.

\subsubsection{4.2 Results}\label{results-1}

Using physical lepton masses (m\_e = 0.51100, m\_mu = 105.658, m\_tau =
1776.86, m\_p = 926.5 MeV (GWT: 6\emph{pi\^{}5}m\_e)):

{\def\LTcaptype{none} % do not increment counter
\begin{longtable}[]{@{}
  >{\raggedright\arraybackslash}p{(\linewidth - 8\tabcolsep) * \real{0.1129}}
  >{\raggedright\arraybackslash}p{(\linewidth - 8\tabcolsep) * \real{0.1774}}
  >{\raggedright\arraybackslash}p{(\linewidth - 8\tabcolsep) * \real{0.2742}}
  >{\raggedright\arraybackslash}p{(\linewidth - 8\tabcolsep) * \real{0.2097}}
  >{\raggedright\arraybackslash}p{(\linewidth - 8\tabcolsep) * \real{0.2258}}@{}}
\toprule\noalign{}
\begin{minipage}[b]{\linewidth}\raggedright
Angle
\end{minipage} & \begin{minipage}[b]{\linewidth}\raggedright
Predicted
\end{minipage} & \begin{minipage}[b]{\linewidth}\raggedright
NuFIT 6.0 (NO)
\end{minipage} & \begin{minipage}[b]{\linewidth}\raggedright
Uncertainty
\end{minipage} & \begin{minipage}[b]{\linewidth}\raggedright
Pull (sigma)
\end{minipage} \\
\midrule\noalign{}
\endhead
\bottomrule\noalign{}
\endlastfoot
theta\_12 (solar) & 33.49 degrees & 33.41 degrees & 0.78 degrees &
+0.1 \\
theta\_23 (atmospheric) & 49.28 degrees & 49.20 degrees & 1.05 degrees &
+0.1 \\
theta\_13 (reactor) & 8.63 degrees & 8.57 degrees & 0.12 degrees &
+0.5 \\
\end{longtable}
}

Mean pull: 0.2 sigma. All three angles within 1 sigma.

The PMNS CP phase is the bare tetrahedral dihedral angle:

\begin{verbatim}
delta_PMNS = arccos(-1/3) = 109.47 degrees
\end{verbatim}

The experimental value is poorly constrained: approximately 180 to 270
degrees (NuFIT 6.0). The prediction is consistent with current data.

\subsubsection{4.3 Structural comparison: CKM
vs.~PMNS}\label{structural-comparison-ckm-vs.-pmns}

{\def\LTcaptype{none} % do not increment counter
\begin{longtable}[]{@{}
  >{\raggedright\arraybackslash}p{(\linewidth - 4\tabcolsep) * \real{0.2368}}
  >{\raggedright\arraybackslash}p{(\linewidth - 4\tabcolsep) * \real{0.3421}}
  >{\raggedright\arraybackslash}p{(\linewidth - 4\tabcolsep) * \real{0.4211}}@{}}
\toprule\noalign{}
\begin{minipage}[b]{\linewidth}\raggedright
Feature
\end{minipage} & \begin{minipage}[b]{\linewidth}\raggedright
CKM (quarks)
\end{minipage} & \begin{minipage}[b]{\linewidth}\raggedright
PMNS (leptons)
\end{minipage} \\
\midrule\noalign{}
\endhead
\bottomrule\noalign{}
\endlastfoot
Power law & 1/2 (surface) & 1/3 (bulk) \\
Base matrix & Identity (small mixing) & TBM (large mixing) \\
Correction & Mass ratios set all angles & Single rotation corrects
TBM \\
CP phase & arccos(5/12) = 65.4 degrees & arccos(-1/3) = 109.5 degrees \\
Confinement & All quarks inside proton & Leptons free on lattice \\
\end{longtable}
}

The CKM angles are small because quark mass ratios m\_light/m\_heavy are
small --- the three generation modes are nearly orthogonal. The PMNS
angles are large because the neutrino sector starts from a democratic
(TBM) base --- the three modes couple nearly equally. The charged-lepton
correction is small (10 degrees) because m\_e/m\_mu is small.

\begin{center}\rule{0.5\linewidth}{0.5pt}\end{center}

\subsection{5. Discussion}\label{discussion}

\subsubsection{5.1 What is derived vs.~what is
assumed}\label{what-is-derived-vs.-what-is-assumed}

The formulas presented here require the following inputs:

\begin{itemize}
\tightlist
\item
  \textbf{d = 3} spatial dimensions (determines the 1/2 and 1/3 powers,
  the tetrahedral angle, and the boundary correction).
\item
  \textbf{Fermion masses} (either GWT-predicted or experimentally
  measured).
\end{itemize}

They do not require:

\begin{itemize}
\tightlist
\item
  Any fitted parameters.
\item
  Any choice of ansatz beyond the PDG parametrization (CKM) and TBM base
  (PMNS).
\item
  Any assumptions about the underlying dynamics beyond ``quarks couple
  at surfaces, leptons couple in bulk.''
\end{itemize}

The surface/bulk distinction is the single structural choice. It is
motivated by the confinement of quarks inside hadrons versus the free
propagation of leptons, which is an observed physical fact, not a model
assumption.

\subsubsection{5.2 Relation to existing
work}\label{relation-to-existing-work}

The Gatto-Sartori-Tonin relation V\_us \textasciitilde{} sqrt(m\_d/m\_s)
and its quadrature extension are well-established {[}1, 2{]}. The
identification V\_cb = sqrt(m\_u/m\_c) appears to be new to this work.
The TBM correction by a single rotation has been explored in various
forms {[}6, 7{]}, but the specific axis formula using
(m\_tau/m\_p)\^{}(1/3) and the connection to lattice wrapping is new.

The tetrahedral origin of CP phases has been discussed in discrete
symmetry models {[}8{]}, but the specific formula cos(delta\_CKM) = 5/12
with its boundary correction interpretation appears to be new.

\subsubsection{5.3 Predictions and tests}\label{predictions-and-tests}

The framework makes several testable predictions:

\begin{enumerate}
\def\labelenumi{\arabic{enumi}.}
\item
  \textbf{PMNS CP phase}: delta\_PMNS = 109.5 degrees. Current data from
  T2K and NOvA are consistent but imprecise. DUNE and Hyper-Kamiokande
  will measure this to approximately 5-10 degrees precision within the
  next decade.
\item
  \textbf{Sensitivity to mass measurements}: Because all mixing angles
  are algebraic functions of mass ratios, improved quark mass
  measurements will sharpen the predictions. The largest current
  uncertainty comes from V\_ub, which depends on sqrt(m\_u/m\_t). The up
  quark mass has approximately 30\% experimental uncertainty, which
  propagates to approximately 15\% uncertainty in V\_ub.
\item
  \textbf{No new parameters}: If additional mixing parameters are ever
  measured (e.g., sterile neutrino mixing), the framework predicts they
  must also be expressible as mass ratios raised to integer-reciprocal
  powers of d.
\end{enumerate}

\subsubsection{5.4 Limitations}\label{limitations}

The geometric origin of the formulas is motivated by the lattice
framework (Section 2) but not rigorously derived from first principles.
The connection between standing-wave overlap integrals and the specific
mass-ratio powers (1/2, 1/3) is plausible but has not been proven at the
level of mathematical rigor expected of a fundamental theory. Similarly,
the boundary correction to the CKM CP phase (1/3 to 5/12) has a clear
geometric interpretation but awaits a formal derivation.

This paper intentionally presents the formulas as empirical relations
that work, alongside a geometric framework that motivates them. The fact
that eight parameters spanning six orders of magnitude are reproduced
from zero free parameters is, in our view, sufficient to merit attention
regardless of the theoretical framework's completeness.

\begin{center}\rule{0.5\linewidth}{0.5pt}\end{center}

\subsection{6. Summary}\label{summary}

Eight mixing parameters of the Standard Model --- four CKM and four PMNS
--- are reproduced from closed-form expressions involving only fermion
mass ratios and the spatial dimension d = 3. The CKM matrix achieves
0.64\% mean accuracy across all nine elements (maximum pull 1.4 sigma).
The PMNS angles are all within 1 sigma of observation. Both CP phases
derive from the tetrahedral dihedral angle in three dimensions. No
parameters are fitted.

The structural distinction between the two sectors --- square root
(surface) for quarks, cube root (bulk) for leptons --- reflects the
observed physics of confinement. These results suggest that the flavor
structure of the Standard Model may have a geometric origin in the
dimensionality of space.

\begin{center}\rule{0.5\linewidth}{0.5pt}\end{center}

\subsection{References}\label{references}

{[}1{]} R. Gatto, G. Sartori, M. Tonin, ``Weak self-masses, Cabibbo
angle, and broken SU(2) x SU(2),'' Phys. Lett. B 28, 128 (1968).

{[}2{]} H. Fritzsch, Z. Xing, ``Mass and flavor mixing schemes of quarks
and leptons,'' Prog. Part. Nucl. Phys. 45, 1 (2000).

{[}3{]} P.F. Harrison, D.H. Perkins, W.G. Scott, ``Tri-bimaximal mixing
and the neutrino oscillation data,'' Phys. Lett. B 530, 167 (2002).

{[}4{]} J.D. Wollenberg, ``Geometric Wave Theory: Fermion Mass Spectrum
from Lattice Breathers,'' Zenodo (2026). {[}To appear{]}

{[}5{]} H. Fritzsch, ``Calculating the Cabibbo angle,'' Phys. Lett. B
70, 436 (1977).

{[}6{]} S.F. King, C. Luhn, ``Neutrino mass and mixing with discrete
symmetry,'' Rep.~Prog. Phys. 76, 056201 (2013).

{[}7{]} G. Altarelli, F. Feruglio, ``Discrete flavor symmetries and
models of neutrino mixing,'' Rev.~Mod. Phys. 82, 2701 (2010).

{[}8{]} I. de Medeiros Varzielas, G.G. Ross, ``Discrete family symmetry,
Higgs mediators and theta\_13,'' JHEP 1212, 041 (2012).

\begin{center}\rule{0.5\linewidth}{0.5pt}\end{center}

\emph{Jonathan D. Wollenberg --- ORCID: 0009-0009-5872-9076}
\emph{GitHub: https://github.com/S-t-u-r-m/geometric-wave-theory}

\end{document}
